% Pursuing Stacks (cotes 134-2 to 134-8): archive cote 134-7, site page 42 (PDF page 43), typescript page 484.
% typescript. Draft compared: ps6.tex:2038-2124.
% AI-transcribed by Claude Opus 5.5 (claude-opus-5-5, adaptive thinking, effort high), Message Batches API, batch ps3 (msgbatch_01XdFmuVVBAcgc7EFxdtCqsj), 2026-11-12; not yet checked by eye.
\scanpage{134-7/42}
\mspage{484}
course,  this Module L(X)  would be unique up to unique isomorphism,  as
the solution of a universal problem embodied by \unsure{(4)},  namely as the
quasi-coherent Module representing the functor
\[ (7) \qquad F \longmapsto \mathrm{Hom}(X,F) \simeq \Gamma(X,F_X) \]
on the category this time of all \uline{quasicoherent} $\underline{O}_k$-Modules.

For the unipotent schematization story,  we are more specifically
interested in the case when X comes from a quasicoherent Module itself,
by forgetting it's module structure.  Now \xout{the} as well known,  the
functor
\[ (8) \qquad M \mapsto W(M) \quad : \quad \mathrm{Ab}_k = (k\text{-Mod}) \qquad \begin{array}[t]{l} \text{category of quasicoherent} \\ \underline{O}_k\text{-Modules, } \hand{\mathrm{Qucoh}(k) \text{ say}} \end{array} \]
is an equivalence of categories.  Thus for X defined by such an M,  the
question of representability of (7)  within the category of quasicoherent
Modules,  amounts to the similar question in $\mathrm{Ab}_k$  for the functor
\[ (7') \qquad N \longmapsto \mathrm{Hom}(W(M),W(N)) \qquad , \]
where the Hom denotes homomorphisms of sheaves \uline{of sets} of course.  Now,
as suggested first\hand{,} somewhat vaguely still\hand{,} in section 111 (page 450),
we have an alternative expression of this functor\hand{,} \xout{as} \hand{via}
\[ (9) \qquad \mathrm{Hom}(W(M),W(N)) \simeq \mathrm{Homcont}_k(\Gamma_k^{\wedge}(M),N) \simeq \xout{\varinjlim \Gamma_k(M)(i} \]
\margin{\textcircled{$\mathrm{Hom}_k($}} \hand{\unsure{$\nearrow$|caret}}
\[ \varinjlim_{i} \mathrm{Hom}\xout{\mathrm{cont}}_k(\Gamma_k(M)(i),N) \quad , \]
where in the second member $\mathrm{Homcont}_k$ denotes the set of k-homomorphisms
\inserted[typed]{\xout{product}}
which are continuous \xout{for the topology induces by the product of} on
\[ (10) \qquad \Gamma_k^{\wedge}(M) = \hand{\prod_{i \geqslant 0}} \Gamma_k^{\,i}(M) \]
\unsure{(}endowed with the product of discrete topologies\unsure{)},  and in the third
we have written
\[ (10') \qquad \Gamma_k(M)(i) = \prod_{j \leqslant i} \Gamma_k^{\,j}(M) \]
for the product of the i first factors occuring in (10).   The map (9)
is \xout{induced by} deduced in evident way from the exponential map
\[ (11) \qquad M \longrightarrow W^{\wedge}\Gamma_k^{\wedge}(M) \overset{\mathrm{def}}{=} \varprojlim_{i} W(\Gamma_k(M)(i)) \qquad . \]
(NB.  The relation between the description (9) of maps $W(M) \rightarrow W(N)$ with
the description given p.451  in terms of maps (p) from $\Gamma_k^{\wedge}(M)$ to
$\Gamma_k^{\wedge}(N)$ ,  is by associating to such a map f it's composition with the
projection of the target upon it's factor N ...).  An uncorrect way of
expressing (9),  which I slipped into in section 111 and kind of remained
in till now,  is by pretending that the k-module  $\Gamma_k^{\wedge}(M)$ represents the

\margin{42}
